\documentclass[11pt,a4paper]{article}

\usepackage[margin=2.5cm]{geometry}
\usepackage{amsmath,amssymb,amsfonts}
\usepackage{booktabs}
\usepackage{array}
\usepackage{longtable}
\usepackage{hyperref}
\usepackage{xcolor}

\title{BeamLib Theory Manual \\ Chapter 2: Euler--Bernoulli 3D Beam Element}
\author{BeamLib Development Team}
\date{Phase 1, Batch 3}

\newcommand{\dd}{\mathrm{d}}
\newcommand{\bmath}[1]{\boldsymbol{#1}}

\begin{document}
\maketitle

\section{Scope and Assumptions}

This chapter derives the 2-node spatial Euler--Bernoulli (EB) beam element used in BeamLib for slender 3D members ($L/h > 20$). It extends the planar EB2D element of Chapter 1 by adding torsion, out-of-plane bending, and a 3D coordinate transformation.

\paragraph{Kinematic assumptions.}
\begin{enumerate}
    \item Plane sections remain plane and perpendicular to the deformed beam axis (no shear deformation, defining EB assumption).
    \item Axial extension along the local $x$-axis is uncoupled from bending and torsion.
    \item Torsion is St.\ Venant uniform torsion: a single rotation DOF $\theta_x$ per node, no warping.
    \item Displacements and rotations are infinitesimal (linear theory).
    \item The two bending planes ($xy$ and $xz$) are uncoupled in the local frame; coupling in the global frame arises only through the coordinate transformation.
    \item Material is linear elastic, isotropic, homogeneous within the section.
\end{enumerate}

\section{Sign and DOF Conventions}

\subsection{Local frame}\label{sec:local-frame}

For each element, BeamLib constructs an orthonormal right-handed local frame $(\hat{x}_\ell, \hat{y}_\ell, \hat{z}_\ell)$ from the two end positions $\bmath{x}_A$, $\bmath{x}_B$ and a per-element \emph{reference vector} $\bmath{r}_{\text{ref}}$:
\begin{align*}
\hat{x}_\ell &= (\bmath{x}_B - \bmath{x}_A) / \|\bmath{x}_B - \bmath{x}_A\|, \\
\hat{y}_\ell &= \mathrm{normalise}\!\bigl(\bmath{r}_{\text{ref}} - (\bmath{r}_{\text{ref}}\!\cdot\!\hat{x}_\ell)\,\hat{x}_\ell\bigr), \\
\hat{z}_\ell &= \hat{x}_\ell \times \hat{y}_\ell.
\end{align*}
The local $y$-axis is the Gram--Schmidt projection of $\bmath{r}_{\text{ref}}$ onto the plane normal to the beam axis. The default $\bmath{r}_{\text{ref}} = (0,1,0)$ matches the BeamLib default in \texttt{ElementConn}; users override $\bmath{r}_{\text{ref}}$ per element to orient the cross-section.

\paragraph{Fallback ladder.} If $\bmath{r}_{\text{ref}}\,\|\,\hat{x}_\ell$ (the projection is the zero vector), BeamMath3D falls back to $\bmath{r}_{\text{ref}} = (0,0,1)$. If that is also parallel to the beam axis (the beam lies along $\hat{z}$ and the user supplied $\bmath{r}_{\text{ref}} = (0,0,1)$), it falls back further to $\bmath{r}_{\text{ref}} = (0,1,0)$. The two fallbacks together cover every beam orientation: $(0,0,1)$ and $(0,1,0)$ span a plane that is non-parallel to any axis, so at least one of them must produce a nonzero projection. The fallback order is fixed (not configurable) so the constructed frame is deterministic.

\subsection{Per-node DOFs}

Each node carries six DOFs in the order
\[
\bmath{d}_i = \begin{bmatrix} u_{x,i} & u_{y,i} & u_{z,i} & \theta_{x,i} & \theta_{y,i} & \theta_{z,i} \end{bmatrix}^\mathsf{T}.
\]
\textbf{Translations $u_x, u_y, u_z$} are along the corresponding global axes (the \emph{global} coordinates the user supplies for nodes / loads). In the canonical local frame in which the element integrates ($\hat{x}_\ell, \hat{y}_\ell, \hat{z}_\ell$), the same indices are interpreted as the local components $u_{x,\ell}, u_{y,\ell}, u_{z,\ell}$.

\textbf{Rotations $\theta_x, \theta_y, \theta_z$} are stored as a 3-component pseudovector. In the local frame:
\begin{itemize}
    \item $\theta_x^\ell$ is the right-hand-rule torsion angle about $+\hat{x}_\ell$.
    \item $\theta_y^\ell$ follows the BeamLib \emph{structural} convention: $\theta_y^\ell = \dd u_z^\ell/\dd x^\ell$. This matches EB2D (Chapter 1, \S2). It is opposite in sign to a right-hand-rule rotation about $+\hat{y}_\ell$.
    \item $\theta_z^\ell$ also follows the structural convention: $\theta_z^\ell = \dd u_y^\ell/\dd x^\ell$. For the $xy$-plane this happens to be the same as the right-hand-rule rotation about $+\hat{z}_\ell$ (because $R_z^{\mathrm{RH}}(\theta)\hat{x}_\ell = \cos\theta\,\hat{x}_\ell + \sin\theta\,\hat{y}_\ell$).
\end{itemize}
The global rotation components are produced by transforming the local pseudovector with $\bmath{\lambda}^\mathsf{T}$, where $\bmath{\lambda}$ is the proper-orthogonal matrix described in \S\ref{sec:transformation}. For elements aligned with the global axes and $\bmath{r}_{\text{ref}} = (0,1,0)$, $\bmath{\lambda} = \bmath{I}$, so the user-facing global rotations directly equal the local structural rotations.

\subsection{Element DOF vector}

A 2-node element has 12 DOFs ordered $[\bmath{d}_A; \bmath{d}_B]$:
\[
\bmath{d}^e = \begin{bmatrix} u_{xA} & u_{yA} & u_{zA} & \theta_{xA} & \theta_{yA} & \theta_{zA} & u_{xB} & u_{yB} & u_{zB} & \theta_{xB} & \theta_{yB} & \theta_{zB} \end{bmatrix}^\mathsf{T}.
\]
C++ indices run 0--11.

\subsection{Sign convention table}

\begin{table}[h!]
\centering
\small
\begin{tabular}{l c l}
\toprule
Quantity & Local index & Positive direction \\
\midrule
Axial displacement       & $u_x$         & along $+\hat{x}_\ell$ (beam axis) \\
Transverse displacement  & $u_y$         & along $+\hat{y}_\ell$ \\
Transverse displacement  & $u_z$         & along $+\hat{z}_\ell$ \\
Torsion rotation         & $\theta_x$    & right-hand rule about $+\hat{x}_\ell$ \\
$xz$-plane rotation      & $\theta_y$    & structural: $\theta_y = \dd u_z/\dd x$ (opposite sign vs.\ RH about $+\hat{y}_\ell$) \\
$xy$-plane rotation      & $\theta_z$    & structural: $\theta_z = \dd u_y/\dd x$ (same as RH about $+\hat{z}_\ell$) \\
Axial force on node      & $F_x$         & along $+\hat{x}_\ell$ \\
Transverse force         & $F_y, F_z$    & along $+\hat{y}_\ell, +\hat{z}_\ell$ \\
Torque                   & $T_x$         & work-conjugate to $\theta_x$ (RH about $+\hat{x}_\ell$) \\
Bending moment           & $M_y$         & work-conjugate to $\theta_y$ (structural convention) \\
Bending moment           & $M_z$         & work-conjugate to $\theta_z$ (structural convention) \\
\bottomrule
\end{tabular}
\caption{EB3D local-frame sign conventions.}
\end{table}

\section{Strain Components}

\subsection{Axial strain}

With axial displacement $u_x(x)$ along the beam axis,
\[
\varepsilon_{x} = \frac{\dd u_x}{\dd x},
\]
and the axial strain energy per unit length is $\tfrac{1}{2} EA\,(\dd u_x/\dd x)^2$.

\subsection{Torsional strain}

For St.\ Venant uniform torsion with rotation $\theta_x(x)$ about the beam axis, the cross-section rotates rigidly and the shear strain in the cross-section produces the strain energy density
\[
\frac{1}{2}\,G I_x \!\left(\frac{\dd \theta_x}{\dd x}\right)^2
\]
per unit length, where $I_x$ is the polar second moment of area (\texttt{SectionProperties::Ix}).

\subsection{Bending in the $xz$ plane}

Deflection $u_z(x)$, rotation $\theta_y(x)$. EB kinematics: $\theta_y = \dd u_z/\dd x$, curvature $\kappa_y = \dd^2 u_z/\dd x^2$. Strain energy density:
\[
\frac{1}{2}\,E I_y \!\left(\frac{\dd^2 u_z}{\dd x^2}\right)^2,
\]
where $I_y = \int_A z^2\,\dd A$ is stored in \texttt{SectionProperties::Iy}. This is the direct 3D analogue of EB2D's xz-plane bending; the only difference is that the symbol is $I_y$ rather than $I_z$ because in 3D the moment of inertia is named after the bending axis ($\hat{y}_\ell$), not after the transverse direction.

\subsection{Bending in the $xy$ plane}

Deflection $u_y(x)$, rotation $\theta_z(x)$. EB kinematics: $\theta_z = \dd u_y/\dd x$, curvature $\kappa_z = \dd^2 u_y/\dd x^2$. Strain energy density:
\[
\frac{1}{2}\,E I_z \!\left(\frac{\dd^2 u_y}{\dd x^2}\right)^2,
\]
with $I_z = \int_A y^2\,\dd A$ stored in \texttt{SectionProperties::Iz}.

\subsection{Total strain energy}

The four contributions are decoupled in the local frame:
\[
U^e = \int_0^L \!\left[\tfrac{1}{2}\,EA\,u_x'^2 + \tfrac{1}{2}\,GI_x\,\theta_x'^2 + \tfrac{1}{2}\,EI_y\,(u_z'')^2 + \tfrac{1}{2}\,EI_z\,(u_y'')^2\right] \dd x.
\]
Coupling between bending and the other modes (axial, torsion) appears only through the transformation $\bmath{T}$ when the local frame is not aligned with the global axes.

\section{Shape Functions}

\begin{itemize}
    \item Axial ($u_x$) and torsion ($\theta_x$): linear, $N_1 = 1 - \xi$, $N_2 = \xi$ on $\xi = x/L \in [0,1]$.
    \item Bending ($u_y$ with $\theta_z = \dd u_y/\dd x$; $u_z$ with $\theta_y = \dd u_z/\dd x$): cubic Hermite, same four basis functions $H_1,\dots,H_4$ as in EB2D (Chapter 1, \S4):
    \[
    H_1 = 1 - 3\xi^2 + 2\xi^3,\quad H_2 = L(\xi - 2\xi^2 + \xi^3),\quad H_3 = 3\xi^2 - 2\xi^3,\quad H_4 = L(-\xi^2 + \xi^3).
    \]
    These reproduce the standard $4\times 4$ Hermite bending stiffness $\frac{EI}{L^3}\,[12, 6L, -12, 6L;\; 6L, 4L^2, -6L, 2L^2;\; -12, -6L, 12, -6L;\; 6L, 2L^2, -6L, 4L^2]$ in DOF order $(u, \theta, u, \theta)$.
\end{itemize}

\section{Weak Form (Local Frame)}\label{sec:weak-form}

The principle of virtual work in the local frame separates into four independent variational statements (one per uncoupled mode) because the four strain contributions are decoupled in $U^e$ (\S3.5):
\begin{align*}
\delta W^e_{\text{axial}}  &= \int_0^L EA\,\delta u_x'\,u_x'\,\dd x  - \bigl[F_x\,\delta u_x\bigr]_A^B, \\
\delta W^e_{\text{tors}}   &= \int_0^L G I_x\,\delta \theta_x'\,\theta_x'\,\dd x - \bigl[T_x\,\delta \theta_x\bigr]_A^B, \\
\delta W^e_{xz\text{ bend}}&= \int_0^L E I_y\,\delta u_z''\,u_z''\,\dd x - \bigl[F_z\,\delta u_z + M_y\,\delta\theta_y\bigr]_A^B, \\
\delta W^e_{xy\text{ bend}}&= \int_0^L E I_z\,\delta u_y''\,u_y''\,\dd x - \bigl[F_y\,\delta u_y + M_z\,\delta\theta_z\bigr]_A^B.
\end{align*}
Discretising with the shape functions of \S4 yields the local element stiffness $\bmath{k}_l \in \mathbb{R}^{12 \times 12}$ whose nonzero blocks are listed in \S\ref{sec:stiffness}.

\section{Local Element Stiffness Matrix}\label{sec:stiffness}

In the canonical local frame (node $A$ at origin, node $B$ at $(L, 0, 0)$), the symmetric local stiffness $\bmath{k}_l$ has four block structures, each restricted to a small subset of the 12 DOFs.

\subsection{Axial block (indices $0, 6$)}
\[
k_l(0,0) = k_l(6,6) = \frac{EA}{L}, \qquad k_l(0,6) = k_l(6,0) = -\frac{EA}{L}.
\]

\subsection{Torsion block (indices $3, 9$)}
\[
k_l(3,3) = k_l(9,9) = \frac{G I_x}{L}, \qquad k_l(3,9) = k_l(9,3) = -\frac{G I_x}{L}.
\]

\subsection{$xy$-plane bending block (indices $1, 5, 7, 11$, uses $E I_z$)}
DOF subset $(u_{yA}, \theta_{zA}, u_{yB}, \theta_{zB})$. The standard Hermite block is
\[
\bmath{k}^{xy}_b = \frac{E I_z}{L^3}\begin{bmatrix}
12 & 6L & -12 & 6L \\
6L & 4L^2 & -6L & 2L^2 \\
-12 & -6L & 12 & -6L \\
6L & 2L^2 & -6L & 4L^2
\end{bmatrix}.
\]
Embedded into the $12\times 12$ matrix at the four indices $(1, 5, 7, 11)$:
\begin{align*}
k_l(1,1) &= k_l(7,7) = \tfrac{12 E I_z}{L^3},   & k_l(1,7)  &= k_l(7,1)  = -\tfrac{12 E I_z}{L^3}, \\
k_l(1,5) &= k_l(5,1) = \tfrac{6 E I_z}{L^2},    & k_l(1,11) &= k_l(11,1) =  \tfrac{6 E I_z}{L^2}, \\
k_l(5,7) &= k_l(7,5) = -\tfrac{6 E I_z}{L^2},   & k_l(7,11) &= k_l(11,7) = -\tfrac{6 E I_z}{L^2}, \\
k_l(5,5) &= k_l(11,11) = \tfrac{4 E I_z}{L},    & k_l(5,11) &= k_l(11,5) =  \tfrac{2 E I_z}{L}.
\end{align*}

\subsection{$xz$-plane bending block (indices $2, 4, 8, 10$, uses $E I_y$)}
DOF subset $(u_{zA}, \theta_{yA}, u_{zB}, \theta_{yB})$. Identical Hermite form with $E I_y$ in place of $E I_z$:
\begin{align*}
k_l(2,2) &= k_l(8,8) = \tfrac{12 E I_y}{L^3},   & k_l(2,8)  &= k_l(8,2)  = -\tfrac{12 E I_y}{L^3}, \\
k_l(2,4) &= k_l(4,2) = \tfrac{6 E I_y}{L^2},    & k_l(2,10) &= k_l(10,2) =  \tfrac{6 E I_y}{L^2}, \\
k_l(4,8) &= k_l(8,4) = -\tfrac{6 E I_y}{L^2},   & k_l(8,10) &= k_l(10,8) = -\tfrac{6 E I_y}{L^2}, \\
k_l(4,4) &= k_l(10,10) = \tfrac{4 E I_y}{L},    & k_l(4,10) &= k_l(10,4) =  \tfrac{2 E I_y}{L}.
\end{align*}
\textbf{Note on the sign pattern.} Even though $\theta_y$ uses the structural convention (opposite RH) while $\theta_z$ matches RH about $+\hat{z}_\ell$, the algebraic Hermite block has the same form in both planes: the $\theta$ DOFs are defined as $\dd u/\dd x$ by construction, so the sign pattern of the coupling entries (positive on the $(u, \theta)$ entries at the same node, negative on the cross-node $(\theta, u)$ entries) is identical in both planes.

\subsection{Residual convention}

For a linear element the local residual is
\[
\bmath{r}_l = \bmath{k}_l\,\bmath{d}_l \in \mathbb{R}^{12}.
\]
At the global level the Newton residual is $\bmath{R} = \bmath{R}_{\mathrm{int}} - \bmath{F}_{\mathrm{ext}}$ and the correction equation is $\bmath{K}\,\delta = -\bmath{R}$, matching the EB2D / linear-solver convention used throughout BeamLib.

\section{Local Consistent Mass Matrix}\label{sec:mass}

The consistent mass $\bmath{m}_l \in \mathbb{R}^{12 \times 12}$ also decomposes into four uncoupled blocks. Let $\rho A L \equiv m_T$ and $\rho I_x L \equiv m_R$.

\subsection{Axial translational mass (indices $0, 6$)}
\[
m_l(0,0) = m_l(6,6) = \tfrac{2}{6}\,m_T, \qquad m_l(0,6) = m_l(6,0) = \tfrac{1}{6}\,m_T.
\]

\subsection{Torsional rotary inertia (indices $3, 9$)}
\[
m_l(3,3) = m_l(9,9) = \tfrac{2}{6}\,m_R, \qquad m_l(3,9) = m_l(9,3) = \tfrac{1}{6}\,m_R.
\]
\textbf{This is the only rotary-inertia term we keep in EB3D.} It comes from integrating $\tfrac{1}{2}\rho I_x \dot\theta_x^2$ over the element with linear shape functions for $\theta_x$ and is \emph{mandatory}: without it the $\theta_x$ DOF carries no mass, the polar inertia is absent, and any torsional eigenmode is missing or shows up as a spurious zero-frequency mode. PROJECT\_SPEC \S3.4 documents this requirement, and Codex Round~3 issue~2 explicitly flagged that omitting this term would invalidate Batch~8/9 torsional modal analysis.

\subsection{Bending blocks (indices $1, 5, 7, 11$ and $2, 4, 8, 10$)}

The standard $4\times 4$ Hermite consistent mass block (in DOF order $(u, \theta, u, \theta)$) is
\[
\bmath{m}_b = \frac{\rho A L}{420}\begin{bmatrix}
156 & 22L & 54 & -13L \\
22L & 4L^2 & 13L & -3L^2 \\
54  & 13L & 156 & -22L \\
-13L & -3L^2 & -22L & 4L^2
\end{bmatrix}.
\]
The same block is placed at indices $(1, 5, 7, 11)$ for $xy$-plane bending and at indices $(2, 4, 8, 10)$ for $xz$-plane bending. Both blocks scale with $\rho A L$ (\emph{not} $\rho I_y$ or $\rho I_z$): the Hermite mass arises from translational kinetic energy $\tfrac{1}{2}\rho A \dot u^2$, where the deflection is interpolated using the four Hermite functions involving $u$ and $\theta$.

\paragraph{No bending rotary inertia.} BeamLib EB3D does \emph{not} add a $\rho I_y \dot\theta_y^2$ or $\rho I_z \dot\theta_z^2$ term, consistent with PROJECT\_SPEC \S3.4 and standard textbook EB mass matrices. The diagonal entries $m_l(5,5) = m_l(11,11) = 4 L^2 \rho A L / 420$ etc.\ are \emph{not} rotary inertia; they are the Hermite-shape-function self-coupling of the translational mass. Timoshenko (Batches~4/5) will add explicit rotary inertia.

\section{Coordinate Transformation}\label{sec:transformation}

\subsection{Per-node block}

Let $\bmath{\lambda} \in \mathbb{R}^{3 \times 3}$ be the matrix whose rows are the local basis vectors expressed in global coordinates:
\[
\bmath{\lambda} = \begin{bmatrix} \hat{x}_\ell^\mathsf{T} \\ \hat{y}_\ell^\mathsf{T} \\ \hat{z}_\ell^\mathsf{T} \end{bmatrix}.
\]
$\bmath{\lambda}$ is proper orthogonal ($\bmath{\lambda}^\mathsf{T}\bmath{\lambda} = \bmath{I}$, $\det\bmath{\lambda} = +1$) by construction (\S\ref{sec:local-frame}). For any 3-vector $\bmath{v}$, $\bmath{v}_\ell = \bmath{\lambda}\,\bmath{v}_g$.

\subsection{Element-level transformation}\label{sec:element-transformation}

The $12\times 12$ element transformation $\bmath{T}$ has a \emph{two-flavour} block-diagonal structure: translations use the plain orthogonal $\bmath{\lambda}$, but rotations use a sign-adapted form $\bmath{S}\bmath{\lambda}\bmath{S}$, where $\bmath{S} = \mathrm{diag}(1,-1,1)$.

\paragraph{Why the rotation block needs $\bmath{S}\bmath{\lambda}\bmath{S}$ instead of $\bmath{\lambda}$.}\label{sec:slambdaSlambdaS}
The BeamLib rotation storage convention is mixed (\S2.2):
\[
\theta_x = +\Omega_x, \qquad \theta_y = -\Omega_y, \qquad \theta_z = +\Omega_z,
\]
where $\bmath{\Omega}$ is the right-hand-rule physical rotation pseudovector. Equivalently, in any orthonormal frame,
\[
\bmath{\theta} = \bmath{S}\,\bmath{\Omega}, \qquad \bmath{S} = \mathrm{diag}(1,-1,1), \qquad \bmath{S}^2 = \bmath{I}.
\]
Because $\bmath{\Omega}$ is a true pseudovector, it transforms by $\bmath{\Omega}_\ell = \bmath{\lambda}\,\bmath{\Omega}_g$ under a proper-orthogonal change of frame. The BeamLib-stored rotation triplet therefore transforms as
\[
\bmath{\theta}_\ell = \bmath{S}\,\bmath{\Omega}_\ell = \bmath{S}\,\bmath{\lambda}\,\bmath{\Omega}_g = \bmath{S}\,\bmath{\lambda}\,\bmath{S}\,(\bmath{S}\,\bmath{\Omega}_g) = \bigl(\bmath{S}\,\bmath{\lambda}\,\bmath{S}\bigr)\,\bmath{\theta}_g.
\]
So the correct global-to-local mapping for the stored rotation DOFs is $\bmath{S}\bmath{\lambda}\bmath{S}$, not $\bmath{\lambda}$. The plain $\bmath{\lambda}$ block is only correct when $\bmath{\lambda}$ commutes with $\bmath{S}$, i.e.\ when $\lambda_{01}=\lambda_{10}=\lambda_{12}=\lambda_{21}=0$ (beams whose local frame keeps the $y$-axis aligned with the global $y$-axis, e.g.\ beams in the $xz$-plane with $\bmath{r}_{\text{ref}}=(0,1,0)$). For beams whose $\bmath{\lambda}$ mixes the $y$-axis with another axis, the commutator $[\bmath{\lambda},\bmath{S}]\ne\bmath{0}$ and using $\bmath{\lambda}$ produces a non-physical local rotation field and a non-zero residual under a true rigid-body motion.

\paragraph{Element transformation matrix.}
With $\bmath{\lambda}_{\text{rot}} \equiv \bmath{S}\bmath{\lambda}\bmath{S}$, the element transformation is
\[
\bmath{T} = \begin{bmatrix}
\bmath{\lambda} & & & \\
& \bmath{\lambda}_{\text{rot}} & & \\
& & \bmath{\lambda} & \\
& & & \bmath{\lambda}_{\text{rot}}
\end{bmatrix},
\]
with the four diagonal blocks corresponding to (node-A translations, node-A rotations, node-B translations, node-B rotations). Element-wise, $\bmath{\lambda}_{\text{rot}}(i,j) = s_i\,\bmath{\lambda}(i,j)\,s_j$ with $\bmath{s} = (1,-1,1)$, equivalently a row-1 \emph{and} column-1 sign flip on $\bmath{\lambda}$.

\paragraph{Internal consistency.}
$\bmath{T}$ remains orthogonal because each diagonal block is orthogonal: $\bmath{\lambda}^\mathsf{T}\bmath{\lambda} = \bmath{I}$ and $(\bmath{S}\bmath{\lambda}\bmath{S})^\mathsf{T}(\bmath{S}\bmath{\lambda}\bmath{S}) = \bmath{S}\bmath{\lambda}^\mathsf{T}\bmath{S}\bmath{S}\bmath{\lambda}\bmath{S} = \bmath{S}\bmath{\lambda}^\mathsf{T}\bmath{\lambda}\bmath{S} = \bmath{S}^2 = \bmath{I}$. The strain-energy invariance and the rigid-body patch test both follow.

\paragraph{Note on torsion ($\theta_x$).}
Even though $\theta_x = \Omega_x$ (no structural sign flip on the $x$ component), the $\bmath{S}\bmath{\lambda}\bmath{S}$ adapter is still correct for the $\theta_x$ row/column: with $s_0 = +1$, the entries involving $\theta_x$ in $\bmath{\lambda}_{\text{rot}}$ pick up factors of $+1 \cdot s_j$. The net effect is that the $\theta_x$ component still satisfies $\theta_{x,\ell} = (\bmath{\lambda}\,\bmath{\Omega}_g)_x$ when expressed back in terms of $\bmath{\Omega}_g$ (verified by direct expansion of $\bmath{S}\bmath{\lambda}\bmath{S}\,\bmath{\theta}_g = \bmath{S}\bmath{\lambda}\bmath{\Omega}_g$).

\subsection{Application convention}

BeamLib applies $\bmath{T}$ as
\[
\bmath{d}_\ell = \bmath{T}\,\bmath{d}_g, \qquad
\bmath{r}_g = \bmath{T}^\mathsf{T}\bmath{r}_\ell, \qquad
\bmath{k}_g = \bmath{T}^\mathsf{T}\bmath{k}_\ell\,\bmath{T}, \qquad
\bmath{m}_g = \bmath{T}^\mathsf{T}\bmath{m}_\ell\,\bmath{T}.
\]
This is the same convention used by EB2D (Chapter 1, \S7) and by \texttt{BeamModel::assemble} / \texttt{BeamModel::assembleMass}.

\subsection{Sanity checks}

\begin{itemize}
    \item Horizontal beam along $+\hat{x}$ with $\bmath{r}_{\text{ref}} = (0,1,0)$: $\hat{x}_\ell = (1,0,0)$, $\hat{y}_\ell = (0,1,0)$, $\hat{z}_\ell = (0,0,1)$, $\bmath{\lambda} = \bmath{I}_3$. Local and global frames coincide; matches EB2D in the $xz$-plane bending sub-block (\texttt{test\_eb2d\_eb3d\_crosscheck.cpp}).
    \item Beam along $+\hat{y}$ with $\bmath{r}_{\text{ref}} = (0,1,0)$ (parallel to beam): first fallback to $(0,0,1)$ is used. Resulting $\bmath{\lambda} = [(0,1,0); (0,0,1); (1,0,0)]$ (\texttt{test\_eb3d\_transform.cpp} subtest C1 verifies this).
    \item Beam along $+\hat{z}$ with $\bmath{r}_{\text{ref}} = (0,0,1)$ (parallel): \emph{both} the original refVector and the first fallback are parallel; second fallback to $(0,1,0)$ produces $\bmath{\lambda} = [(0,0,1); (0,1,0); (-1,0,0)]$ (\texttt{test\_eb3d\_transform.cpp} subtest C2).
    \item Orthogonality and energy invariance follow from $\bmath{\lambda}^\mathsf{T}\bmath{\lambda} = \bmath{I}_3$ in each $3\times 3$ block.
\end{itemize}

\section{Global Assembly}

For each element with connectivity $(A, B)$ and reference vector $\bmath{r}_{\text{ref}}^e$:
\begin{enumerate}
    \item Build $\bmath{T}^e$ from $\bmath{x}_A$, $\bmath{x}_B$, $\bmath{r}_{\text{ref}}^e$ via \texttt{BeamMath3D::buildTransformation3D}.
    \item Gather the 12-vector $\bmath{d}_g^e$ from nodal storage. Fixed DOFs carry their prescribed value (via \texttt{Node::totalDof}), matching the Batch 2B convention for nonzero prescribed displacement.
    \item $\bmath{d}_\ell^e = \bmath{T}^e\,\bmath{d}_g^e$.
    \item Pass $\bmath{x}_A^{\text{loc}} = (0,0,0)$, $\bmath{x}_B^{\text{loc}} = (L,0,0)$ ($L = \|\bmath{x}_B - \bmath{x}_A\|$), and $\bmath{d}_\ell^e$ to \texttt{EulerBernoulli3D::computeElement}. The local stiffness depends only on $L$ and the section properties.
    \item $\bmath{r}_g^e = (\bmath{T}^e)^\mathsf{T}\bmath{r}_\ell^e$, $\bmath{k}_g^e = (\bmath{T}^e)^\mathsf{T}\bmath{k}_\ell^e\,\bmath{T}^e$.
    \item Scatter $\bmath{r}_g^e$, $\bmath{k}_g^e$ into the free-DOF system using \texttt{DofMap}.
\end{enumerate}
The same flow applies to the mass matrix using \texttt{computeMass}.

\section{Reaction Recovery and Internal Forces}

Reactions follow the Batch 2B mechanism unchanged: full-system reassembly of $\bmath{R}_{\mathrm{int}}^{\text{full}}$, then $R_d^{\mathrm{reaction}} = R_d^{\mathrm{int,full}} - F_d^{\mathrm{ext}}$ at every fixed DOF. The torsion test (\texttt{test\_eb3d\_torsion.cpp}) checks that the recovered $\theta_x$ reaction equals $-T_x$ when a tip torque $T_x$ is applied: at equilibrium the internal residual at the root $\theta_x$ DOF is $-T_x$ (the elements push back against the applied torque), and the external load there is zero, so the reaction is $-T_x$.

\paragraph{Cross-section internal forces.} The 3D set of internal forces is $\{N, V_y, V_z, T, M_y, M_z\}$, six quantities per cross-section. The Phase~1 \texttt{ElementInternalForces} struct (used by EB2D) currently exposes only the 2D subset $\{N, V_z, M_y\}$. Extending it to 3D is deferred to a later batch; the EB3D \texttt{computeInternalForces} stub populates only $\{N, V_z, M_y\}$ using the same sign convention as EB2D and explicitly notes that $V_y$, $M_z$, and $T$ are out of scope until the struct is extended. No EB3D test in Batch~3 reads internal forces; the torsion reaction test goes through \texttt{computeReactions}, which uses the full element residual rather than these fields.

\section{Formula-to-C++ Index Mapping}

\begin{longtable}{c c l}
\toprule
Formula symbol & C++ index & Meaning \\
\midrule \endhead
$u_{xA}$        & \texttt{dispVec[0]}  & node-A local axial translation \\
$u_{yA}$        & \texttt{dispVec[1]}  & node-A local $y$ translation \\
$u_{zA}$        & \texttt{dispVec[2]}  & node-A local $z$ translation \\
$\theta_{xA}$   & \texttt{dispVec[3]}  & node-A torsion rotation \\
$\theta_{yA}$   & \texttt{dispVec[4]}  & node-A $xz$-plane rotation (structural) \\
$\theta_{zA}$   & \texttt{dispVec[5]}  & node-A $xy$-plane rotation (structural) \\
$u_{xB}$        & \texttt{dispVec[6]}  & node-B local axial translation \\
$u_{yB}$        & \texttt{dispVec[7]}  & node-B local $y$ translation \\
$u_{zB}$        & \texttt{dispVec[8]}  & node-B local $z$ translation \\
$\theta_{xB}$   & \texttt{dispVec[9]}  & node-B torsion rotation \\
$\theta_{yB}$   & \texttt{dispVec[10]} & node-B $xz$-plane rotation (structural) \\
$\theta_{zB}$   & \texttt{dispVec[11]} & node-B $xy$-plane rotation (structural) \\
\midrule
$k_l(0,0)$, $k_l(6,6)$    & \texttt{ke(0,0)}, \texttt{ke(6,6)}   & $+EA/L$ \\
$k_l(0,6)$                & \texttt{ke(0,6)}                     & $-EA/L$ \\
$k_l(3,3)$, $k_l(9,9)$    & \texttt{ke(3,3)}, \texttt{ke(9,9)}   & $+G I_x/L$ \\
$k_l(3,9)$                & \texttt{ke(3,9)}                     & $-G I_x/L$ \\
$k_l(1,1)$, $k_l(7,7)$    & \texttt{ke(1,1)}, \texttt{ke(7,7)}   & $+12 E I_z/L^3$ (xy bending) \\
$k_l(1,7)$                & \texttt{ke(1,7)}                     & $-12 E I_z/L^3$ \\
$k_l(1,5)$, $k_l(1,11)$   & \texttt{ke(1,5)}, \texttt{ke(1,11)}  & $+6 E I_z/L^2$ \\
$k_l(5,7)$, $k_l(7,11)$   & \texttt{ke(5,7)}, \texttt{ke(7,11)}  & $-6 E I_z/L^2$ \\
$k_l(5,5)$, $k_l(11,11)$  & \texttt{ke(5,5)}, \texttt{ke(11,11)} & $+4 E I_z/L$ \\
$k_l(5,11)$               & \texttt{ke(5,11)}                    & $+2 E I_z/L$ \\
$k_l(2,2)$, $k_l(8,8)$    & \texttt{ke(2,2)}, \texttt{ke(8,8)}   & $+12 E I_y/L^3$ (xz bending) \\
$k_l(2,8)$                & \texttt{ke(2,8)}                     & $-12 E I_y/L^3$ \\
$k_l(2,4)$, $k_l(2,10)$   & \texttt{ke(2,4)}, \texttt{ke(2,10)}  & $+6 E I_y/L^2$ \\
$k_l(4,8)$, $k_l(8,10)$   & \texttt{ke(4,8)}, \texttt{ke(8,10)}  & $-6 E I_y/L^2$ \\
$k_l(4,4)$, $k_l(10,10)$  & \texttt{ke(4,4)}, \texttt{ke(10,10)} & $+4 E I_y/L$ \\
$k_l(4,10)$               & \texttt{ke(4,10)}                    & $+2 E I_y/L$ \\
\midrule
$m_l(0,0)$, $m_l(6,6)$    & \texttt{me(0,0)}, \texttt{me(6,6)}   & $+\tfrac{2}{6}\rho A L$ \\
$m_l(0,6)$                & \texttt{me(0,6)}                     & $+\tfrac{1}{6}\rho A L$ \\
$m_l(3,3)$, $m_l(9,9)$    & \texttt{me(3,3)}, \texttt{me(9,9)}   & $+\tfrac{2}{6}\rho I_x L$ \\
$m_l(3,9)$                & \texttt{me(3,9)}                     & $+\tfrac{1}{6}\rho I_x L$ \\
$m_l(1,1)$, $m_l(7,7)$    & \texttt{me(1,1)}, \texttt{me(7,7)}   & $+156 \rho A L/420$ (xy bending) \\
$m_l(2,2)$, $m_l(8,8)$    & \texttt{me(2,2)}, \texttt{me(8,8)}   & $+156 \rho A L/420$ (xz bending) \\
(others as listed in \S\ref{sec:mass}) & & \\
\midrule
$\hat{x}_\ell$            & \texttt{LocalFrame3D::Vx}            & beam axis $(\bmath{x}_B - \bmath{x}_A)/L$ \\
$\hat{y}_\ell$            & \texttt{LocalFrame3D::Vy}            & projection of \texttt{refVector} normal to $\hat{x}_\ell$ \\
$\hat{z}_\ell$            & \texttt{LocalFrame3D::Vz}            & $\hat{x}_\ell \times \hat{y}_\ell$ \\
$\bmath{\lambda}$         & \texttt{LocalFrame3D::lambda}        & rows = $(\hat{x}_\ell^\mathsf{T}, \hat{y}_\ell^\mathsf{T}, \hat{z}_\ell^\mathsf{T})$ \\
$\bmath{T}$               & \texttt{buildTransformation3D}        & translation blocks use $\bmath{\lambda}$; rotation blocks use $\bmath{S}\bmath{\lambda}\bmath{S}$ (\S\ref{sec:slambdaSlambdaS}) \\
$\bmath{S}$               & in-line in code                       & $\mathrm{diag}(1,-1,1)$ structural-vs-RH sign adapter for $\theta_y$ \\
\bottomrule
\caption{EB3D formula-to-C++ index mapping.}
\end{longtable}

\section{Verification: Analytical Cases}\label{sec:verification}

\subsection{Cantilever tip transverse load $F_z$}

Beam along $+\hat{x}$, fixed at node 0, tip force $F_z$ in global $z$. With $\bmath{r}_{\text{ref}} = (0,1,0)$, $\bmath{\lambda} = \bmath{I}_3$, so the problem reduces to local $xz$-plane bending with stiffness $E I_y$:
\[
u_z(L) = \frac{F_z L^3}{3 E I_y}, \qquad \theta_y(L) = \frac{F_z L^2}{2 E I_y}.
\]
For $F_z = -1000$, $E = 200{\times}10^9$, $I_y = 8.333{\times}10^{-6}$, $L = 1.0$: $u_z(L) = -2{\times}10^{-4}$, $\theta_y(L) = -3{\times}10^{-4}$. Both are negative for downward load (\texttt{test\_eb3d\_cantilever.cpp} case 1). All unrelated DOFs (i.e., $u_x$, $u_y$, $\theta_x$, $\theta_z$) are zero by construction because no other modes are excited.

\subsection{Cantilever tip transverse load $F_y$}

Same beam, tip force $F_y$ in global $y$. With $\bmath{\lambda} = \bmath{I}_3$, the $xy$-plane bending engages with stiffness $E I_z$:
\[
u_y(L) = \frac{F_y L^3}{3 E I_z}, \qquad \theta_z(L) = \frac{F_y L^2}{2 E I_z}.
\]
Both negative for $F_y < 0$ (\texttt{test\_eb3d\_cantilever.cpp} case 2).

\subsection{Cantilever tip torque}

Same beam, tip torque $T_x$ about $+\hat{x}$. Pure St.\ Venant torsion gives a linearly varying $\theta_x$:
\[
\theta_x(L) = \frac{T_x L}{G I_x},
\]
with all other DOFs zero. Root reaction in $\theta_x$ equals $-T_x$ (\texttt{test\_eb3d\_torsion.cpp}).

\subsection{Cross-check: EB3D vs.\ EB2D in the $xz$ plane}

For a beam aligned with $+\hat{x}$ under tip load $F_z$, EB2D ($I_z$ for $xz$ bending) and EB3D ($I_y$ for $xz$ bending) must agree exactly when $I_y^{3D} = I_z^{2D}$. The cross-check test compares tip $u_z$ and $\theta_y$ and confirms relative agreement below $10^{-9}$ (\texttt{test\_eb2d\_eb3d\_crosscheck.cpp}). The test sets the unused $I_z^{3D}$ to a different value as a guard against silent mis-indexing of the two bending planes.

\subsection{Rigid-body patch test (transformation correctness)}

\texttt{test\_eb3d\_rigid\_body.cpp} imposes a small physical rigid-body motion $\bmath{u} = \bmath{t} + \bmath{\Omega}\times\bmath{X}$ with $\bmath{\theta} = \bmath{S}\bmath{\Omega}$ on a 2-node element with arbitrary orientation and $\bmath{r}_{\text{ref}}$. The element internal residual norm must be at machine zero (rigid-body $\Rightarrow$ no strain $\Rightarrow$ no internal force).

The geometry is deliberately chosen so $\bmath{\lambda}$ does not commute with $\bmath{S}$ (a beam in the $xy$-plane plus a generic $\bmath{r}_{\text{ref}}$): if the rotation transformation used the plain $\bmath{\lambda}$ block instead of $\bmath{S}\bmath{\lambda}\bmath{S}$, the test would produce a residual of $\mathcal{O}(\|EI\bmath{\Omega}/L\|)$ rather than zero. This is the test Codex Batch~3 review issue~1 requested, and it directly exercises \S\ref{sec:slambdaSlambdaS}.

\subsection{Coordinate-transformation tests}

Three subtests in \texttt{test\_eb3d\_transform.cpp}:
\begin{enumerate}
    \item \textbf{Local transverse force on a rotated beam:} 45-degree element in the $xz$ plane with $\bmath{r}_{\text{ref}} = (0,1,0)$. A local-$z$ force is transformed into global coordinates, the model is solved, and the tip displacement is transformed back to the local frame. The local-frame result equals the horizontal-beam reference solution to $10^{-9}$.
    \item \textbf{Reference vector changes the bending plane:} a horizontal beam with $\bmath{r}_{\text{ref}} = (0,1,0)$ versus $\bmath{r}_{\text{ref}} = (0,0,1)$ under the same global $F_y$. With $I_y \ne I_z$ the two tip $u_y$ values differ; each matches the analytical formula for the corresponding bending plane.
    \item \textbf{Fallback ladder:} beam along $+\hat{y}$ with $\bmath{r}_{\text{ref}} = (0,1,0)$ (first fallback) and beam along $+\hat{z}$ with $\bmath{r}_{\text{ref}} = (0,0,1)$ (second fallback). In both cases the constructed $\bmath{\lambda}$ is orthogonal with $\det = +1$, and a transverse global force produces the analytical EB deflection along the bending plane induced by the fallback.
\end{enumerate}

\section{Verification Tolerance Table}

\begin{table}[h!]
\centering
\small
\begin{tabular}{l l}
\toprule
Check & Tolerance \\
\midrule
Cantilever tip $u_z$, $\theta_y$ under $F_z$ (10 elements) & relative $10^{-9}$ \\
Cantilever tip $u_y$, $\theta_z$ under $F_y$ (10 elements) & relative $10^{-9}$ \\
Cantilever tip $\theta_x$ under $T_x$ (10 elements)       & relative $10^{-9}$ \\
Linearly-varying $\theta_x$ along span                   & relative $10^{-9}$ \\
Torsion root reaction $\theta_x$ = $-T_x$                 & relative $10^{-9}$ \\
Unrelated DOF amplitudes (e.g., $u_x$ under $F_z$)        & absolute $10^{-14}$ \\
Reaction at unloaded fixed DOF                            & absolute $10^{-9}$ \\
EB3D vs.\ EB2D cross-check ($u_z$, $\theta_y$)            & relative $10^{-9}$ \\
Rigid-body residual norm (arbitrary orientation)          & absolute $10^{-9}$ \\
Mass matrix entries (local, single horizontal element)    & relative $10^{-12}$ \\
Mass matrix entries (transformed, generic orientation)    & relative $10^{-10}$ \\
Mass matrix symmetry                                       & absolute $10^{-12}$ \\
Rotated-beam tip displacement (local frame round-trip)    & relative $10^{-9}$ \\
Fallback frame orthogonality ($\|\bmath{\lambda}\bmath{\lambda}^\mathsf{T} - \bmath{I}\|$) & absolute $10^{-12}$ \\
Fallback-frame tip deflection under transverse load        & relative $10^{-9}$ \\
NR iteration count (linear case)                          & exactly 1 correction step \\
Final NR residual (linear case)                           & $\le 10^{-7}$ absolute (NRConfig.tol) \\
\bottomrule
\end{tabular}
\caption{Verification tolerances for the EB3D Batch 3 test suite.}
\end{table}

\section{Implementation Pitfalls (for the Reviewer)}

\begin{enumerate}
    \item \textbf{Section property choice per plane.} EB3D $xy$-plane bending uses $\texttt{props.Iz}$; $xz$-plane bending uses $\texttt{props.Iy}$. Swapping them is a silent bug. The cross-check test explicitly sets the unused $I$ to a different value to guard this.
    \item \textbf{$\bmath{r}_{\text{ref}}$ is per-element.} \texttt{ElementConn::refVector} defaults to $(0,1,0)$ but the user may set it differently per element. The local frame depends only on $\bmath{r}_{\text{ref}}$ and the beam axis, never on absolute nodal positions (in contrast to the reference DEM-FEM C++ implementation, which uses $\texttt{refNode} - \bmath{x}_A$ and is not translation-invariant).
    \item \textbf{Fallback determinism.} The fallback ladder is $(0,0,1) \to (0,1,0)$, fixed regardless of which $\bmath{r}_{\text{ref}}$ the user originally supplied. Tests assert specific frame components for both fallbacks.
    \item \textbf{Torsional rotary inertia is mandatory.} The polar inertia term $\rho I_x$ in the $\theta_x$ block is required for any modal or implicit-dynamic analysis. PROJECT\_SPEC \S3.4 and Codex Round~3 issue~2 both highlight this.
    \item \textbf{Bending rotary inertia is intentionally absent.} Diagonal entries on $\theta_y$, $\theta_z$ in the mass matrix come from translational kinetic energy through the Hermite shape functions, not from $\rho I$. Timoshenko (Batches~4/5) adds explicit rotary inertia.
    \item \textbf{$\theta_y$ structural vs.\ right-hand-rule convention.} EB3D inherits EB2D's structural convention $\theta_y = \dd u_z / \dd x$, which is opposite in sign to the right-hand-rule rotation about $+\hat{y}_\ell$. The numerical Hermite block algebra is unchanged (the structural definition matches what the Hermite shape function $H_2(\xi)$ enforces), but the \emph{rotation transformation block} must be $\bmath{S}\bmath{\lambda}\bmath{S}$ rather than $\bmath{\lambda}$ to be consistent with this convention (\S\ref{sec:slambdaSlambdaS}). Using plain $\bmath{\lambda}$ silently passes any test where $\bmath{\lambda}$ commutes with $\bmath{S}$ (e.g.\ beams in the $xz$-plane with $\bmath{r}_{\text{ref}}=(0,1,0)$) but fails the 3D rigid-body patch test for general orientations.
    \item \textbf{Internal force postprocessing is partial.} \texttt{computeInternalForces} currently exposes only $\{N, V_z, M_y\}$; $V_y$, $M_z$, $T$ require an extension to the shared \texttt{ElementInternalForces} struct and are deferred to a later batch.
\end{enumerate}

\end{document}
